% GENERATED FILE. Edit canonical lessons, then run npm run generate:books.
% canonical-source-hash: fnv1a64:418a467dffe364a2
% canonical-lessons: JA-C91-yamu, JA-C91-shimau, JA-C91-hirou, JA-C91-tetsudau, JA-C91-mayou, JA-R91-first-pass-this-and-that-and-everyday-actions, JA-R91-second-pass-more-everyday-actions

\chapter{To stop, To put away, To pick up, To help, To get lost}
\label{ch:ja-to-stop-to-put-away-to-pick-up-to-help-to-get-lost}
\hlchaptermodality{91}

\begin{chapteropening}
\textbf{By the end of this chapter:} \emph{I can say to stop, to put away, to pick up, to help and to get lost.}

Complete the last lesson of chapter 91: i can say to stop, to put away, to pick up, to help and to get lost.
\end{chapteropening}

\section[\texorpdfstring{\ja{やむ} (yamu)}{やむ (yamu)}]{\ja{やむ} (yamu) — to stop (rain)}

\label{lesson:JA-C91-yamu}



\begin{quote}
Before the new one: say the Japanese for to go regularly, then the Japanese for to wrap.
\end{quote}

\subsection*{You'll want to know: \ja{やむ}}
\textbf{\ja{やむ}} — \emph{yamu} — \textquotedblleft{}to stop (rain)\textquotedblright{}.

\begin{grammarlens}[title={What you've built}]
One more word for this chapter.
\end{grammarlens}

\subsection*{Your turn}
\begin{itemize}
\raggedright
  \item \emph{Say it:} \emph{yamu}
  \item \emph{Say it:} \emph{yamu}, once more
  \item \emph{Say it:} say \emph{yamu}
  \item \emph{Recall:} say the Japanese for to go regularly, then the Japanese for to wrap, then say \emph{yamu} again
\end{itemize}

\begin{tcolorbox}[breakable,colback=teal!4,colframe=teal!35!black,title={Before you move on}]
What does \ja{やむ} mean? (\textquotedblleft{}To stop (rain)\textquotedblright{}.) Say it once more.
\end{tcolorbox}
\section[\texorpdfstring{\ja{しまう} (shimau)}{しまう (shimau)}]{\ja{しまう} (shimau) — to put away}

\label{lesson:JA-C91-shimau}



\begin{quote}
Before the new one: say the Japanese for to wrap, then the Japanese for to stop.
\end{quote}

\subsection*{You'll want to know: \ja{しまう}}
\textbf{\ja{しまう}} — \emph{shimau} — \textquotedblleft{}to put away\textquotedblright{}.

\begin{grammarlens}[title={What you've built}]
One more word for this chapter.
\end{grammarlens}

\subsection*{Your turn}
\begin{itemize}
\raggedright
  \item \emph{Say it:} \emph{shimau}
  \item \emph{Say it:} \emph{shimau}, once more
  \item \emph{Say it:} say \emph{shimau}
  \item \emph{Recall:} say the Japanese for to wrap, then the Japanese for to stop, then say \emph{shimau} again
\end{itemize}

\begin{tcolorbox}[breakable,colback=teal!4,colframe=teal!35!black,title={Before you move on}]
What does \ja{しまう} mean? (\textquotedblleft{}To put away\textquotedblright{}.) Say it once more.
\end{tcolorbox}
\section[\texorpdfstring{\ja{ひろう} (hirou)}{ひろう (hirou)}]{\ja{ひろう} (hirou) — to pick up}

\label{lesson:JA-C91-hirou}



\begin{quote}
Before the new one: say the Japanese for to stop, then the Japanese for to put away.
\end{quote}

\subsection*{You'll want to know: \ja{ひろう}}
\textbf{\ja{ひろう}} — \emph{hirou} — \textquotedblleft{}to pick up\textquotedblright{}.

\begin{grammarlens}[title={What you've built}]
One more word for this chapter.
\end{grammarlens}

\subsection*{Your turn}
\begin{itemize}
\raggedright
  \item \emph{Say it:} \emph{hirou}
  \item \emph{Say it:} \emph{hirou}, once more
  \item \emph{Say it:} say \emph{hirou}
  \item \emph{Recall:} say the Japanese for to stop, then the Japanese for to put away, then say \emph{hirou} again
\end{itemize}

\begin{tcolorbox}[breakable,colback=teal!4,colframe=teal!35!black,title={Before you move on}]
What does \ja{ひろう} mean? (\textquotedblleft{}To pick up\textquotedblright{}.) Say it once more.
\end{tcolorbox}
\section[\texorpdfstring{\ja{てつだう} (tetsudau)}{てつだう (tetsudau)}]{\ja{てつだう} (tetsudau) — to help}

\label{lesson:JA-C91-tetsudau}



\begin{quote}
Before the new one: say the Japanese for to put away, then the Japanese for to pick up.
\end{quote}

\subsection*{You'll want to know: \ja{てつだう}}
\textbf{\ja{てつだう}} — \emph{tetsudau} — \textquotedblleft{}to help\textquotedblright{}.

\begin{grammarlens}[title={What you've built}]
One more word for this chapter.
\end{grammarlens}

\subsection*{Your turn}
\begin{itemize}
\raggedright
  \item \emph{Say it:} \emph{tetsudau}
  \item \emph{Say it:} \emph{tetsudau}, once more
  \item \emph{Say it:} say \emph{tetsudau}
  \item \emph{Recall:} say the Japanese for to put away, then the Japanese for to pick up, then say \emph{tetsudau} again
\end{itemize}

\begin{tcolorbox}[breakable,colback=teal!4,colframe=teal!35!black,title={Before you move on}]
What does \ja{てつだう} mean? (\textquotedblleft{}To help\textquotedblright{}.) Say it once more.
\end{tcolorbox}
\section[\texorpdfstring{\ja{まよう} (mayou)}{まよう (mayou)}]{\ja{まよう} (mayou) — to get lost}

\label{lesson:JA-C91-mayou}



\begin{quote}
Before the new one: say the Japanese for to pick up, then the Japanese for to help.
\end{quote}

\subsection*{You'll want to know: \ja{まよう}}
\textbf{\ja{まよう}} — \emph{mayou} — \textquotedblleft{}to get lost\textquotedblright{}.

\begin{grammarlens}[title={What you've built}]
One more word for this chapter.
\end{grammarlens}

\subsection*{Your turn}
\begin{itemize}
\raggedright
  \item \emph{Say it:} \emph{mayou}
  \item \emph{Say it:} \emph{mayou}, once more
  \item \emph{Say it:} say \emph{mayou}
  \item \emph{Recall:} say the Japanese for to pick up, then the Japanese for to help, then say \emph{mayou} again
\end{itemize}

\begin{tcolorbox}[breakable,colback=teal!4,colframe=teal!35!black,title={Before you move on}]
What does \ja{まよう} mean? (\textquotedblleft{}To get lost\textquotedblright{}.) Say it once more.
\end{tcolorbox}
\section[(dialogue)]{This and that, and everyday actions — a first pass}

\label{lesson:JA-R91-first-pass-this-and-that-and-everyday-actions}



\begin{quote}
Say the words for to help and to get lost.
\end{quote}

\subsection*{Your turn}
Say these aloud:

\begin{itemize}
\raggedright
  \item a street corner, within the city, the suburbs, the opposite side, upstream
  \item this, that, which, this kind of, that kind of
  \item what kind of, other, everybody, this sort of, that sort of
  \item to hit, to breathe in, to sow; to wind, to cook, to put on
  \item to bloom, to solve, to hold in one's arms, to step on, to bite
\end{itemize}

\begin{tcolorbox}[breakable,colback=teal!4,colframe=teal!35!black,title={Before you move on}]
A street corner? (\textbf{\ja{まちかど}}, \emph{machikado}.) Within the city? (\textbf{\ja{しない}}, \emph{shinai}.) The suburbs? (\textbf{\ja{こうがい}}, \emph{kougai}.) The opposite side? (\textbf{\ja{はんたいがわ}}, \emph{hantaigawa}.) Upstream? (\textbf{\ja{かわかみ}}, \emph{kawakami}.) This? (\textbf{\ja{この}}, \emph{kono}.) That? (\textbf{\ja{あの}}, \emph{ano}.) Which? (\textbf{\ja{どの}}, \emph{dono}.) This kind of? (\textbf{\ja{こんな}}, \emph{konna}.) That kind of? (\textbf{\ja{あんな}}, \emph{anna}.) What kind of? (\textbf{\ja{どんな}}, \emph{donna}.) Other? (\textbf{\ja{ほか}}, \emph{hoka}.) Everybody? (\textbf{\ja{みんな}}, \emph{minna}.) This sort of? (\textbf{\ja{こういう}}, \emph{kouiu}.) That sort of? (\textbf{\ja{ああいう}}, \emph{aaiu}.) To hit? (\textbf{\ja{うつ}}, \emph{utsu}.) To breathe in? (\textbf{\ja{すう}}, \emph{suu}.) To sow; to wind? (\textbf{\ja{まく}}, \emph{maku}.) To cook? (\textbf{\ja{たく}}, \emph{taku}.) To put on? (\textbf{\ja{はく}}, \emph{haku}.) To bloom? (\textbf{\ja{さく}}, \emph{saku}.) To solve? (\textbf{\ja{とく}}, \emph{toku}.) To hold in one's arms? (\textbf{\ja{だく}}, \emph{daku}.) To step on? (\textbf{\ja{ふむ}}, \emph{fumu}.) To bite? (\textbf{\ja{かむ}}, \emph{kamu}.)
\end{tcolorbox}
\section[(dialogue)]{More everyday actions — a second pass}

\label{lesson:JA-R91-second-pass-more-everyday-actions}



\begin{quote}
Say the words for to pile up and to ask for.
\end{quote}

\subsection*{Your turn}
Say these aloud:

\begin{itemize}
\raggedright
  \item to pile up, to ask for, to go forward, to hang out to dry, to point; to hold up
  \item to copy, to try out, to put back, to leave behind, to knock down
  \item to break, to boil, to turn round, to grab, to fold
  \item to be surprised, to arrive, to receive, to go regularly, to wrap
  \item to stop, to put away, to pick up, to help, to get lost
\end{itemize}

\begin{tcolorbox}[breakable,colback=teal!4,colframe=teal!35!black,title={Before you move on}]
To pile up? (\textbf{\ja{つむ}}, \emph{tsumu}.) To ask for? (\textbf{\ja{たのむ}}, \emph{tanomu}.) To go forward? (\textbf{\ja{すすむ}}, \emph{susumu}.) To hang out to dry? (\textbf{\ja{ほす}}, \emph{hosu}.) To point; to hold up? (\textbf{\ja{さす}}, \emph{sasu}.) To copy? (\textbf{\ja{うつす}}, \emph{utsusu}.) To try out? (\textbf{\ja{ためす}}, \emph{tamesu}.) To put back? (\textbf{\ja{もどす}}, \emph{modosu}.) To leave behind? (\textbf{\ja{のこす}}, \emph{nokosu}.) To knock down? (\textbf{\ja{たおす}}, \emph{taosu}.) To break? (\textbf{\ja{こわす}}, \emph{kowasu}.) To boil? (\textbf{\ja{わかす}}, \emph{wakasu}.) To turn round? (\textbf{\ja{まわす}}, \emph{mawasu}.) To grab? (\textbf{\ja{つかむ}}, \emph{tsukamu}.) To fold? (\textbf{\ja{たたむ}}, \emph{tatamu}.) To be surprised? (\textbf{\ja{おどろく}}, \emph{odoroku}.) To arrive? (\textbf{\ja{とどく}}, \emph{todoku}.) To receive? (\textbf{\ja{いただく}}, \emph{itadaku}.) To go regularly? (\textbf{\ja{かよう}}, \emph{kayou}.) To wrap? (\textbf{\ja{つつむ}}, \emph{tsutsumu}.) To stop? (\textbf{\ja{やむ}}, \emph{yamu}.) To put away? (\textbf{\ja{しまう}}, \emph{shimau}.) To pick up? (\textbf{\ja{ひろう}}, \emph{hirou}.) To help? (\textbf{\ja{てつだう}}, \emph{tetsudau}.) To get lost? (\textbf{\ja{まよう}}, \emph{mayou}.)
\end{tcolorbox}
